\label{chap:83-12}

\section{El Núcleo de Transducción Hensel--Witt}
\label{p12-83-c12-s1:chap:nthw}
\index{Núcleo de Transducción Hensel--Witt}

\readerquestion{¿Dónde se cruzan realmente los cilindros, la monodromía, las
cinco regiones, las hexadas, las octadas y los lectores físicos?}

\subsection{Núcleo de Transducción Hensel–Witt: microregiones de \(\pi\), holonomía \(\Gamma_9\), registro dodecafásico y transporte hexada–octada}

El NTHW es el paquete de datos y mapas
\[
\NTHW=
\bigl(
\mathscr R_\pi^{\rm micro},
\Gnine,
\Cdoce,
K,
\mathcal W_{12\to24},
\mathcal I_{6\to8}
\bigr),
\]
con las compatibilidades siguientes:
\begin{enumerate}
  \item las regiones Hensel comparten una palabra cociente y conservan
  distinta memoria;
  \item \(\Gnine\) transporta la orientación y produce monodromía;
  \item \(\Cdoce\) organiza el observable firmado y el sello \(K\);
  \item la lectura Witt convierte soportes del mismo borde en hexadas y
  estrellas;
  \item la extensión \(W_{12}\to W_{24}\) produce la interfaz
  hexada--octada;
  \item los lectores arquimediano, excepcional y dimensional reciben
  invariantes del mismo estado.
\end{enumerate}

\begin{figure}[p]
\centering
\resizebox{\textwidth}{!}{\input{figures/fig_nucleo_hensel_witt}}
\caption{El NTHW como interfaz tipada. La tarea matemática central es hacer
conmutar el diagrama, no acumular coincidencias alrededor de un nombre.}
\label{p12-83-c12-s1:fig:nthw}
\end{figure}

El NTHW da un nombre formal al lugar que antes se describía como «Piedra
Rosetta». La sustitución es conceptual: una piedra traduce por semejanza; una
transducción especifica entrada, salida, dato conservado y pérdida de
información.

\subsection{Del código ternario al sistema de Witt}

La frontera reversible de APP selecciona, con calibres declarados, una matriz
\(A_W\) sobre \(\FF_3\) que satisface
\[
A_WA_W^T=2I\pmod3.
\]
El código gráfico
\[
C_W=\{(q,qA_W):q\in\FF_3^6\}
\]
es el Golay ternario extendido. Su enumerador de pesos es
\[
1+264y^6+440y^9+24y^{12}.
\]
Las \(264\) palabras de peso seis, identificadas por signo, dan \(132\)
soportes. Cada subconjunto de cinco puntos pertenece a una única hexada:
\[
W_{12}=S(5,6,12).
\]

El campo de \(495\) involuciones de estrella genera un grupo de orden
\(95\,040\), reconocido como \(M_{12}\). Una estrella aislada no genera el
grupo. El marco HMT ordenado tiene estabilizador trivial y por ello fija un
calibre dentro de la acción.

\subsection{Hexada, octada y el corredor \texorpdfstring{\(90/120\)}{90/120}}

Sea \(H\) una hexada y \(O_H\) una octada que la contiene. Al marcar un punto
y un par de puntos de la hexada se obtienen
\[
F_6(H)=H\times\binom H2,
\qquad
|F_6|=6\binom62=90.
\]
Si el punto marcado puede recorrer la octada:
\[
F_8(H)=O_H\times\binom H2,
\qquad
|F_8|=8\binom62=120.
\]
La inclusión \(H\hookrightarrow O_H\) induce \(F_6\hookrightarrow F_8\).
Su defecto normalizado es
\[
\boxed{
\frac{90-120}{24\binom62}
=\frac{6-8}{24}
=-\frac1{12}.}
\]
Este teorema combinatorio es anterior a la lectura angular de Barbero. Los
dos objetos comparten la interfaz \(6\to8\), pero no son idénticos.

\subsection{\(4\) rutas hacia -1/12}

La construcción contiene cuatro realizaciones que deben exponerse juntas y no
contarse como cuatro pruebas independientes de una misma procedencia.

\subsubsection{Operador elemental del ciclo dodecafásico y actualización de la fase firmada}

Un diente de \(\Cdoce\) mide \(30^\circ\). Entonces
\[
\frac{30}{120}-\frac{30}{90}
=\frac14-\frac13
=-\frac1{12}.
\]
Es la sombra angular del patrón \(90/120\).

\subsubsection{Matriz de Gram de las secciones supervivientes y criterio de positividad}

En el salto certificado \(60\to81\), dos fibras tienen tamaños \(0\) y
\(12\). El operador de Gram es
\[
K_{60,81}=\operatorname{diag}(0,12).
\]
Sobre el subespacio vivo, su pseudoinverso vale \(1/12\); con la orientación
negativa de borde,
\[
E_{\rm surv}=-K_{60,81}^{+}
\]
produce \(-1/12\).

\subsubsection{Extractor de escala triádico}

La realización mediante un extractor de escala triádico es la construcción por
segunda diferencia desarrollada íntegramente en
\cref{sec:c12-segunda-diferencia-triadica}. Esta primera comparecencia fija su
función dentro del inventario de realizaciones; la residencia indicada conserva
las definiciones de \(T_L\), \(a(L)\), \(A_1(L)\), \(C(L)\), la independencia
del residuo respecto de \(L\) y el estatuto exacto de la normalización.

\subsubsection{Corredor modular entre cartas dodecafásicas y conservación del invariante de fase}

La función eta satisface
\[
\frac{\eta(\tau)}{\eta(3\tau)}
=q^{-1/12}\prod_{n\ge1}(1+q^n+q^{2n})^{-1}.
\]
El exponente procede de \(1/24-3/24=-1/12\). Al tomar doce copias:
\[
t_3(q)=\left(\frac{\eta(\tau)}{\eta(3\tau)}\right)^{12}
=q^{-1}-12+54q-76q^2-243q^3+\cdots.
\]
El coeficiente \(54\) coincide con el defecto APP. El corredor
\[
j=\frac{(t_3+27)(t_3+243)^3}{t_3^3}
\]
conecta después con la rama modular clásica.

\begin{exactbox}[title={Qué se ha demostrado en el corredor}]
Son exactos el defecto APP \(54\), el exponente eta \(-1/12\), el defecto
hexada--octada y las identidades modulares una vez fijado \(t_3\). La
afirmación fuerte es que todos sean imágenes naturales de un único operador
TPK. Esa naturalidad debe demostrarse mediante el diagrama del NTHW; no se
deduce sólo de repetir los mismos enteros.
\end{exactbox}

\subsection{Defecto de potencias primas y operador número}

Para una potencia prima \(p^m\), el defecto completo de APP generaliza a
\[
\Delta_{p^m}^{\rm full}
=\frac{m(p-1)}2p^{2m-1}.
\]
En particular,
\[
\Delta_3=3,
\qquad
\Delta_9=54,
\qquad
\Delta_{27}=729.
\]
Al normalizar,
\[
a_{p,m}
=\frac{2\Delta_{p^m}^{\rm full}}{(p-1)p^{2m-1}}
=m.
\]
Con \(x=p^{-s}\):
\[
\sum_{m\ge1}a_{p,m}x^m
=\sum_{m\ge1}mx^m
=\frac{x}{(1-x)^2}
=x\frac{\dd}{\dd x}\frac1{1-x}.
\]
El factor \((1-x)^{-1}\) es el factor local de Euler; la derivada inserta el
operador número. Así aparece un puente exacto:
\[
\boxed{
\text{torre de defectos APP en }p^m
\longleftrightarrow
\text{inserción del operador número en el factor local de Euler}.}
\]

\subsection{Condiciones de naturalidad conjunta para la transducción arquimediana, excepcional y dimensional}

El NTHW será un teorema completo de transducción cuando exista una familia de
invariantes \(\mathcal I\) y lectores
\[
\mathfrak R,\quad\mathfrak E,\quad\mathfrak D
\]
tales que
\[
\mathcal I
=\mathcal I_{\rm arq}\circ\mathfrak R
=\mathcal I_{\rm exc}\circ\mathfrak E
=\mathcal I_{\rm dim}\circ\mathfrak D,
\]
y cuando los mapas de región, sello, incidencia y acción se deriven del mismo
transporte. Las pruebas finitas ya proporcionan candidatos: monodromía,
soporte \(B_K\), patrón de préstamos, paridad two-way y conteos \(90/120\).

\begin{programbox}[title={Prueba ejecutable 108--144--90--120}]
Etiquetar cada una de las \(1008\) rutas de \(\pi\) por órbita CT108 y por
completación Witt; calcular las fibras del mapa; comprobar si las cuatro
regiones de multiplicidad \(144\) realizan el factor \(8/6\) sobre bloques
de \(108\); y verificar si la región de \(432\) transporta la monodromía. Un
solo contraejemplo a la compatibilidad de etiquetas refutaría este puente en
su forma actual.
\end{programbox}

\begin{takeawaybox}
El NTHW es el centro real del relato porque allí el mismo borde puede leerse
como cilindro, monodromía, sello, incidencia y acción. Su potencia no consiste
en que muchos números se parezcan, sino en que varios mapas finitos ya
comparten soportes y defectos exactos. El trabajo siguiente es probar la
naturalidad conjunta de esas cartas.
\end{takeawaybox}


\section{Inclusión marcada de la incidencia hexada--octada y preservación de sus multiplicidades}

Fijados la dodecada y el alineamiento del
\cref{p12-83-c17-s1:sec:w12-w24-elevacion}, sea \(H\) una hexada en las coordenadas HMT,
sea \(\ell(H)\subset D\) su imagen y sea
\(O_H=\iota_D(\ell(H))\) su elevación única. Al marcar un punto
y un par de puntos de la hexada residual se obtienen
\[
F_6(H)=\ell(H)\times\binom{\ell(H)}2,
\qquad
|F_6|=6\binom62=90.
\]
Si el punto marcado puede recorrer la octada:
\[
F_8(H)=O_H\times\binom{\ell(H)}2,
\qquad
|F_8|=8\binom62=120.
\]
La inclusión \(\ell(H)\hookrightarrow O_H\) induce el morfismo de incidencia
\[
 \mathcal I_{6\to8}:F_6(H)\hookrightarrow F_8(H).
\]
Al normalizar la diferencia por el factor
elegido \(24\binom62\), se obtiene
\[
\boxed{
\frac{90-120}{24\binom62}
=\frac{6-8}{24}
=-\frac1{12}.}
\]
La igualdad es un defecto normalizado de la interfaz combinatoria
\(6\to8\). El factor \(24\binom62\) forma parte de la definición de esta
normalización; la inclusión por sí sola determina la diferencia \(-30\).
El operador \(\mathcal I_{6\to8}\) tiene como dominio configuraciones
marcadas de incidencia; no es el extractor transversal del estado
dodecafásico ni la familia de momentos angulares.

\section{Clasificación por dominio de las \(5\) realizaciones algebraicas de \(-1/12\)}

Además del defecto normalizado de la inclusión hexada--octada, aparecen
cinco realizaciones operatoriales del mismo racional. Sus dominios se
especifican por separado; la igualdad de valores no identifica los
operadores sin un entrelazador adicional.

\subsection{Diferencia de periodicidades dodecafásicas}

Un diente de \(\Cdoce\) mide \(30^\circ\). Entonces
\[
\frac{30}{120}-\frac{30}{90}
=\frac14-\frac13
=-\frac1{12}.
\]
Es la diferencia orientada entre los periodos de \(90^\circ\) y
\(120^\circ\) en esta realización.

\subsection{Pseudoinversa del operador de Gram sobre la componente viva}

En el salto exacto \(60\to81\), dos fibras tienen tamaños \(0\) y
\(12\). El operador de Gram es
\[
K_{60,81}=\operatorname{diag}(0,12).
\]
Sobre la componente no nula de su rango, la pseudoinversa vale \(1/12\); con la orientación
negativa de borde,
\[
E_{\rm surv}=-K_{60,81}^{+}
\]
produce \(-1/12\).

\subsection{2.ª diferencia entre escalas triádicas}
\label{sec:c12-segunda-diferencia-triadica}

Definamos
\[
T_L=\sum_{n=1}^{L-1}n(L-n)=\frac{L(L^2-1)}6,
\qquad
a(L)=-\frac{T_L}{L^2}=-\frac L6+\frac1{6L}.
\]
La primera resta de escala es
\[
A_1(L)=a(L)-\frac{a(3L)}3=\frac4{27L},
\]
y la segunda
\[
C(L)=LA_1(L)-\frac{3L}{3}A_1(3L)=\frac8{81}.
\]
El residuo \(8/81\) es independiente de \(L\). Para obtener
\(-1/12\) se aplica la normalización declarada
\[
R(L)=-\frac{27}{32}C(L)=-\frac1{12}.
\]
El cálculo deriva \(8/81\); no deriva por sí solo el factor \(-27/32\).

\subsection{Composición de transformaciones modulares y preservación del estado firmado}

La función eta satisface
\[
\frac{\eta(\tau)}{\eta(3\tau)}
=q^{-1/12}\prod_{n\ge1}(1+q^n+q^{2n})^{-1}.
\]
El exponente procede de \(1/24-3/24=-1/12\). Al tomar doce copias:
\[
t_3(q)=\left(\frac{\eta(\tau)}{\eta(3\tau)}\right)^{12}
=q^{-1}-12+54q-76q^2-243q^3+\cdots.
\]
El coeficiente \(54\) coincide con el defecto APP\@. La identidad racional
\[
j=\frac{(t_3+27)(t_3+243)^3}{t_3^3}
\]
establece la conexión con la rama modular clásica.

\subsection{Defecto de los proyectores de diferencia cíclica}

Sea \(S\) el desplazamiento cíclico sobre \(\QQ^{12}\) y definamos
\[
 D_r=I-S^r.
\]
El núcleo de \(D_r\) está formado por los vectores constantes en cada órbita
de \(S^r\); por tanto,
\[
 \dim\ker D_r=\gcd(12,r).
\]
Sean \(\Pi_{\ker D_3}\) y \(\Pi_{\ker D_4}\) los proyectores racionales sobre
\(\ker D_3\) y \(\ker D_4\), y sea
\(\tau(T)=\operatorname{Tr}(T)/12\) la traza normalizada. Entonces
\[
 \boxed{
 \tau(\Pi_{\ker D_3}-\Pi_{\ker D_4})
 =\frac{\operatorname{rank}\Pi_{\ker D_3}
       -\operatorname{rank}\Pi_{\ker D_4}}{12}
 =\frac{3-4}{12}
 =-\frac1{12}.}
\]
Esta realización expresa el racional como defecto transversal de los pasos
tres y cuatro en el ciclo de doce fases.

\begin{remark}[Dominios de las realizaciones]
Son exactos el defecto APP \(54\), el exponente eta \(-1/12\), el defecto
combinatorio de la elevación hexada--octada, y las
identidades modulares una vez fijado \(t_3\). También son exactos el valor
de la pseudoinversa, el extractor normalizado y el defecto de proyectores.
Cada igualdad conserva su dominio y su aplicación correspondiente.
\end{remark}
